\section{Behavior Reconstruction}

\rvmt reconstructs behavior in layers. Hardware records provide the event
sequence and architectural context. ELF metadata provides binary identity and
code ranges. Runtime process maps bind the target process and its loaded
objects. Reference logs provide expected behavior for comparison but are not
reported as hardware output.

The local code-attribution fixtures cover cases that commonly break simple
PC-to-binary mapping: ELF SHA-256 identity, PIE and ASLR load bias, dynamic
library ownership, fork/exec ownership, stripped binaries, and source maps not
captured on the board. Source-line information is reported only when it is
bound to the exact board artifact.

This reconstruction model is intentionally limited. It supports evidence such
as syscall pairing, measured-window target attribution, and executable
ownership where the required inputs exist. It does not evaluate a complete OS
process graph, complete filesystem history, complete memory-object tracking,
or recovery of all pointer strings.
